\documentclass[11pt,a4paper]{article}
\usepackage[T1]{fontenc}
\usepackage[margin=1in]{geometry}
\usepackage{amsmath,amssymb,amsthm,mathtools}
\usepackage{booktabs}
\usepackage{xcolor}
\usepackage[colorlinks=true,linkcolor=blue!60!black,citecolor=blue!60!black,urlcolor=blue!60!black]{hyperref}
\usepackage{microtype}
\usepackage{xurl}
\usepackage{enumitem}
\setlength{\emergencystretch}{3em}
\newtheorem{theorem}{Theorem}
\newtheorem{proposition}{Proposition}
\newtheorem{lemma}{Lemma}
\theoremstyle{remark}
\newtheorem{remark}{Remark}
\newcommand{\R}{\mathbb{R}}
\newcommand{\Z}{\mathbb{Z}}
\hypersetup{pdftitle={Complete monotonicity of the image-charge interaction in a grounded slab},pdfauthor={Xavier Callens}}

\title{\textbf{Complete monotonicity of the image-charge interaction\\ in a grounded slab:
a heat-kernel proof and a Lean~4 formalisation}}
\author{Xavier Callens\\ \small SocrateAI Research \\
\small \url{https://github.com/xaviercallens/SocrateAI-Scientific-QuantumFluids}}
\date{October 2026 \\ \small Draft for expert reading; not refereed; written with an AI system}

\begin{document}
\maketitle

\begin{abstract}\noindent
Two pairs of opposite charges (two ``excitons'') interact, inside a slab bounded by two grounded perfect conductors, through the
Dirichlet Green function of the Laplacian. We show that the resulting pair interaction $U(r)$, with $r$ the lateral distance, is
\emph{completely monotone in $r^2$}: $U(r)=\int_0^\infty e^{-r^2\sigma}\mu(\sigma)\,d\sigma$ with a bounded non-negative density
$\mu$. The proof is a one-line consequence of the product structure of the Dirichlet heat kernel of the slab and of the positive
definiteness of its transverse factor; the same argument covers any laterally homogeneous layered dielectric. A Lean~4/Mathlib
formalisation (52 theorems, no \texttt{sorry}, standard axioms) proves the grounded-plane case by a different route, via Poisson
summation. The statement is elementary and may well be known; we ask for that to be checked. We say precisely how it connects to
universal-optimality theorems, and that the planar statement we would need there is not available in a refereed form.
\end{abstract}

\section{Setting and statement}

Let $h_1,h_2>0$, $\ell=h_1+h_2$, $P=2\ell$, and let $\Omega=\R^2\times(-h_2,h_1)$ be the slab between grounded perfect conductors
at $z=h_1$ and $z=-h_2$ (vacuum inside; layered media in \S\ref{sec:diel}). Write $G(x,z;x',z')$ for the Dirichlet Green function of
$-\Delta$ on $\Omega$, and $K=4\pi G$, so that $K=1/R$ in free space. By the method of images, with $R_\pm$ the distances to the
charge and image charges,
\begin{equation}\label{eq:images}
K\big((x,z),(x',z')\big)=\sum_{n\in\Z}\Big[\frac1{\sqrt{r^2+(z-z'-Pn)^2}}-\frac1{\sqrt{r^2+(z+z'-2h_1-Pn)^2}}\Big],
\qquad r=|x-x'|,
\end{equation}
where the series converges absolutely, the two terms of each pair cancelling to $O(n^{-2})$.

An exciton consists of an electron (charge $-1$) in the plane $z=z_e$ and a hole ($+1$) in the plane $z=z_h$ at the same lateral
position; the interaction energy of two excitons at lateral distance $r$, in units of $e^2/\varepsilon$ and without
$r$-independent self-image terms, is
\begin{equation}\label{eq:U}
U(r)=K_{ee}(r)+K_{hh}(r)-2K_{eh}(r),\qquad K_{ab}(r)=K\big((x,z_a),(x',z_b)\big),
\end{equation}
which has $U(r)\sim 2/r$ at $0$ when $z_e\ne z_h$ lie in the open slab (and $U\equiv0$ for $z_e=z_h$). The formula \eqref{eq:images} makes sense
for all real $z,z'$ as the odd periodic extension; the Lean file proves the statement in that generality.

\begin{theorem}\label{thm:main}
For all $z_e,z_h\in[-h_2,h_1]$ and all $h_1,h_2>0$ there is a measurable function $\mu:(0,\infty)\to[0,\infty)$, continuous and bounded, with
\[
U(r)=\int_0^\infty e^{-r^2\sigma}\,\mu(\sigma)\,d\sigma\qquad(r>0).
\]
In particular $t\mapsto U(\sqrt t)$ is completely monotone on $(0,\infty)$: it is $C^\infty$ and $(-1)^k\frac{d^k}{dt^k}U(\sqrt t)\ge0$ for all $k$.
Explicitly, with $\lambda_m=(m\pi/\ell)^2$ and $\varphi_m(z)=\sin\!\big(m\pi(z+h_2)/\ell\big)$,
\begin{equation}\label{eq:mu}
\mu(\sigma)=\frac{4}{P\sigma}\sum_{m\ge1}e^{-\lambda_m/(4\sigma)}\,\big(\varphi_m(z_e)-\varphi_m(z_h)\big)^2\ \ge0 .
\end{equation}
\end{theorem}

\section{Proof via the Dirichlet heat kernel}

\textbf{Step 1 (product structure).} The Dirichlet Laplacian on $\Omega$ is $-\Delta_x\otimes1+1\otimes A$ with $A=-\partial_z^2$ on
$(-h_2,h_1)$ with Dirichlet conditions; the two summands act on different variables and commute, so the heat kernel of $\Omega$ is
$p_s(x,z;x',z')=(4\pi s)^{-1}e^{-|x-x'|^2/4s}\,q_s(z,z')$, where $q_s$ is the kernel of $e^{-sA}$. By the method of images
\[
q_s(z,z')=\sum_{n\in\Z}\big[g_s(z-z'-Pn)-g_s(z+z'-2h_1-Pn)\big],\qquad g_s(u)=(4\pi s)^{-1/2}e^{-u^2/4s},
\]
and also $q_s(z,z')=\frac{2}{\ell}\sum_{m\ge1}e^{-s\lambda_m}\varphi_m(z)\varphi_m(z')$.

\textbf{Step 2 (Laplace representation).} $(-\Delta)^{-1}=\int_0^\infty e^{s\Delta}\,ds$ on $L^2$ with Dirichlet conditions, so
$G=\int_0^\infty p_s\,ds$ and
\[
K(r;z,z')=4\pi G=\int_0^\infty \frac1s\,e^{-r^2/4s}\,q_s(z,z')\,ds=\int_0^\infty e^{-r^2\sigma}\,\frac{q_{1/4\sigma}(z,z')}{\sigma}\,d\sigma
\qquad(\sigma=1/4s).
\]
(The integral converges for $r>0$ because $q_s$ is bounded by $g_s(0)$ and decays like $e^{-s\lambda_1}$ as $s\to\infty$.)
Hence $U(r)=\int_0^\infty e^{-r^2\sigma}\mu(\sigma)\,d\sigma$ with $\mu(\sigma)=\sigma^{-1}S_{1/4\sigma}$ and
\begin{equation}\label{eq:S}
S_s:=q_s(z_e,z_e)+q_s(z_h,z_h)-2\,q_s(z_e,z_h).
\end{equation}

\textbf{Step 3 (positivity, the only point).} By the semigroup property and symmetry of $q$,
$q_s(z,z')=\int q_{s/2}(z,w)\,q_{s/2}(w,z')\,dw$, hence
\[
S_s=\int_{-h_2}^{h_1}\big(q_{s/2}(z_e,w)-q_{s/2}(z_h,w)\big)^2\,dw\ \ge0 .
\]
In the eigenfunction expansion this is $S_s=\frac2\ell\sum_me^{-s\lambda_m}(\varphi_m(z_e)-\varphi_m(z_h))^2$, which with
$s=1/4\sigma$ gives \eqref{eq:mu}. (Numerically, \eqref{eq:images}--\eqref{eq:mu} agree to $10^{-12}$ at $r=0.3,1,2.5$ (symmetric truncation $|n|\le2\times10^5$) for $h_1=1,h_2=1.5,z_e=0.3,z_h=-0.4$: \texttt{python/gated\_note\_check.py}.) (Neutrality of the exciton is what makes \eqref{eq:S} a square: the non-neutral combination
$K_{ee}-2K_{eh}$ can be negative mode by mode; the Lean file contains this as a negative control.)

\textbf{Step 4 (regularity, boundedness).} Take $z_e\ne z_h$. As $\sigma\to0$ every term of \eqref{eq:mu} carries $e^{-\lambda_m/4\sigma}$, so
$\mu(\sigma)\to0$. As $\sigma\to\infty$ ($s\to0$), $S_s=2g_s(0)+O(e^{-c/s})$ because the cross terms $q_s(z_e,z_h)$ and the image terms are
exponentially small in $1/s$ for points at positive distance from each other and from the planes (for points on a plane the image terms
vanish with the Dirichlet condition), hence $\mu(\sigma)\sim 2(\pi\sigma)^{-1/2}$: this is the $2/r$ singularity of $U$ at $0$, and it is
integrable against $e^{-r^2\sigma}$ for every $r>0$. $\mu$ is continuous and bounded on $(0,\infty)$ (the Lean file proves a bound by the
Poisson form for $\sigma\ge1$ and by the first term of the series for $\sigma\le1$). Differentiation under the integral, justified by $\sigma^ke^{-t\sigma}\le C_{k,t_0}e^{-t_0\sigma/2}$ for
$t\ge t_0$, gives $(-1)^k\frac{d^k}{dt^k}\int e^{-t\sigma}\mu=\int\sigma^ke^{-t\sigma}\mu\ge0$ (the easy direction of Bernstein's theorem).
\hfill$\square$

\begin{remark}
Since $e^{-r^2\sigma}$ is a Gaussian in $x\in\R^d$ for every $d$, the representation exhibits $U$ as a non-negative mixture of Gaussians. It is
not a positive-definite \emph{function} ($U(0)=\infty$), but it is a positive-definite distribution on $\R^d$ for $d\ge2$ (its Fourier transform
is positive and finite for $k\ne0$; for $d=1$, $2/|x|$ is not locally integrable). This is Bernstein's easy direction combined with Gaussians,
not Schoenberg's theorem.
\end{remark}

\section{Layered dielectrics (not formalised)}\label{sec:diel}

\begin{proposition}\label{prop:diel}
Let $\varepsilon:(-h_2,h_1)\to[\varepsilon_{\min},\varepsilon_{\max}]\subset(0,\infty)$ be measurable (a stack of homogeneous layers, or a uniaxial pair $\varepsilon_\parallel(z),\varepsilon_\perp(z)$, for which the lateral Laplacian carries $\varepsilon_\parallel/\varepsilon_\perp$ and the argument below is unchanged up to a rescaling of $r$),
with grounded conductors at the two ends. Then \eqref{eq:U} (with $K=4\pi G$, $G$ the Green function of $-\nabla\!\cdot\varepsilon\nabla$) satisfies
$U(r)=\sum_{m\ge1}c_m\,(\psi_m(z_e)-\psi_m(z_h))^2K_0(\sqrt{\lambda_m}\,r)$ with $c_m\ge0$ (in units of $e^2$), where $(\lambda_m,\psi_m)$, $\psi_m$ normalised in $L^2(\varepsilon\,dz)$, are the Dirichlet
eigenpairs of $-\varepsilon^{-1}\partial_z\varepsilon\partial_z$ in $L^2(\varepsilon\,dz)$. Since $K_0(\sqrt{\lambda t})=\tfrac12\int_0^\infty
e^{-ts}e^{-\lambda/4s}\,ds/s$, $U$ is completely monotone in $r^2$.
\end{proposition}

\emph{Proof.} (i) The form $a(u,v)=\int\varepsilon u'v'\,dz$ on $H^1_0(-h_2,h_1)$ is closed, non-negative and has compact resolvent on
$L^2(\varepsilon\,dz)$ (bounded $\varepsilon$, bounded below by $\varepsilon_{\min}$); its operator is $A=-\varepsilon^{-1}\partial_z\varepsilon\partial_z$, with discrete spectrum
$0<\lambda_1\le\lambda_2\le\dots$ and eigenfunctions $\psi_m$ orthonormal in $L^2(\varepsilon\,dz)$. (ii) The potential of a unit charge at $(0,z')$ solves $-\nabla\!\cdot\varepsilon\nabla G=\delta$,
i.e. $(-\Delta_x+A)G=\delta_x\otimes\delta_{z'}/\varepsilon(z')$; since $e^{-sA}$ has a kernel $q_s(z,w)$ with respect to $\varepsilon\,dw$, and
$\int q_s(z,w)\,\delta_{z'}(w)\varepsilon(w)^{-1}\varepsilon(w)\,dw=q_s(z,z')$, we get $G=\int_0^\infty p^{\R^2}_s(x)\,q_s(z,z')\,ds$ exactly as before (the
heat semigroup is a tensor product because the two operators act on different variables). (iii) $q_s(z,z')=\sum_me^{-s\lambda_m}\psi_m(z)\psi_m(z')$ is symmetric
and satisfies $q_s=\int q_{s/2}\,q_{s/2}\,\varepsilon\,dw$, so $S_s=\int(q_{s/2}(z_e,w)-q_{s/2}(z_h,w))^2\varepsilon\,dw=\sum_me^{-s\lambda_m}(\psi_m(z_e)-\psi_m(z_h))^2\ge0$.
(iv) Integrability of $s^{-1}e^{-r^2/4s}S_s$ at $s\to0$ follows from $q_s(z,z)\le Cs^{-1/2}$ (Nash/ultracontractivity in one dimension),
and at $s\to\infty$ from $q_s\le Ce^{-\lambda_1s}$. All terms are non-negative, so by Tonelli the sum and the $s$-integral commute, and
$\int_0^\infty s^{-1}e^{-r^2/4s}e^{-s\lambda}\,ds=2K_0(\sqrt\lambda\,r)$ gives the expansion with $c_m=2$ (the factor $4\pi$ of $K=4\pi G$ cancels the $1/4\pi s$ of the lateral heat kernel); for the grounded
planes with constant $\varepsilon$, $\psi_m=\sqrt{2/\ell}\,\varphi_m$ and this is $U=\frac4\ell\sum_m(\varphi_m(z_e)-\varphi_m(z_h))^2K_0(\sqrt{\lambda_m}\,r)$, which agrees with the image series numerically. \hfill$\square$

\textbf{What this does not cover.} Finite screening by the gates (non-ideal metals), lateral inhomogeneity, and quantum or
polarisation effects are outside the statement; for them the product structure fails.

\section{The Lean~4 formalisation}

\texttt{exploration/exciton/lean/GatedKernel.lean} (1083 lines, 52 theorems; Lean 4.34.1, Mathlib v4.34.1; \texttt{import Mathlib}
only). It proves Theorem~\ref{thm:main} for the grounded planes by the \emph{Fourier} route, not by the heat-kernel argument above:
the image sums are Gaussian lattice sums, Poisson summation turns them into $\sum_me^{-\pi^2m^2/(P^2s)}\cos(\cdot)$, and the combination
\eqref{eq:S} becomes a sum of $S_m=2(1-\cos(b-a))(1+\cos(a+b))=2(\sin a-\sin b)^2$, which is the same quantity as
$(\varphi_m(z_e)-\varphi_m(z_h))^2$ in \eqref{eq:mu}.

\begin{center}\small
\begin{tabular}{@{}p{0.28\textwidth}p{0.66\textwidth}@{}}
\toprule
declaration & content \\ \midrule
\texttt{trig\_core\_sq} & $2-\cos2a-\cos2b-2\cos(b-a)+2\cos(a+b)=2(\sin a-\sin b)^2$ \\
\texttt{gGreen\_top}, \texttt{gGreen\_bottom} & the Gaussian lattice Green function vanishes on both planes (bijective re-indexing of image offsets) \\
\texttt{gsum\_poisson} & (sum over $m\in\Z$) $\sum_ne^{-s(c-Pn)^2}=\frac{\sqrt{\pi/s}}P\sum_me^{-\pi^2m^2/(P^2s)}\cos(2\pi mc/P)$, \emph{derived} from Mathlib's Gaussian/Jacobi-theta results \\
\texttt{Sexc\_nonneg} & $S(s)\ge0$ for all $s>0$, $h_1+h_2>0$, all real $z_e,z_h$ \\
\texttt{muExc\_bounded}, \texttt{muExc\_continuousOn} & $\mu$ is bounded and continuous on $(0,\infty)$ \\
\texttt{lap\_admissible} & Laplace transform of a non-negative density with finite $\int e^{-st}\mu$ is smooth with $(-1)^rg^{(r)}\ge0$ (easy Bernstein) \\
\texttt{Uimg\_eq\_lap}, \texttt{imageSeries\_admissible} & the image series equals the Laplace transform, hence is \texttt{AdmissiblePotential} \\
\texttt{unneutral\_not\_nonneg} & negative control: dropping $V_{hh}$ breaks positivity \\
\bottomrule
\end{tabular}
\end{center}

All 17 \texttt{\#print axioms} outputs are \texttt{[propext, Classical.choice, Quot.sound]}; there is no \texttt{sorry}, \texttt{axiom}, or
\texttt{native\_decide}. The only hypotheses are $0<s,t$ and $0<h_1+h_2$. An independent verification (separate instance) compiled the file and recomputed the
kernel by a direct signed image sum with $N=4000$ (agreement $10^{-9}$ at three radii, and with the Laplace integral to $10^{-12}$).
Definition used: \texttt{AdmissiblePotential g} $:=$ $g$ is $C^\infty$ on $(0,\infty)$, $g\ge0$, and $(-1)^r\,g^{(r)}\ge0$ for all $r$ (a
verbatim copy of the definition in the upstream repository~\cite{OpenAI2026}). It is satisfied by $g=0$, so non-triviality is a separate statement
($U\ge0$ follows from $\mu\ge0$; strict positivity for $z_e\ne z_h$ is not stated in the file). The Mathlib version pinned here has
\texttt{AbsolutelyMonotoneOn} but no completely monotone class and no Laplace transform of measures, so the file carries its own.

\section{Relation to universal optimality}

Cohn and Kumar defined universally optimal configurations on spheres as those minimising the energy for \emph{every} completely monotone function
of squared distance~\cite{CohnKumar2007}, with the Euclidean version as their conjecture; $E_8$ and the Leech lattice were proved universally optimal in dimensions 8 and 24
by Cohn, Kumar, Miller, Radchenko and Viazovska~\cite{CKMRV2022}. The triangular lattice in dimension 2 is conjectured to be
universally optimal in the same sense. A claimed proof of that planar statement is publicly available~\cite{OpenAI2026}, but it is
AI-generated and \textbf{has not been refereed to our knowledge}; we do not use it as a proved fact. What Theorem~\ref{thm:main}
supplies is exactly the hypothesis of that statement:

\begin{proposition}[unconditional, lattices]\label{prop:lat}
Among Bravais lattices $\Lambda\subset\R^2$ of fixed density, the triangular lattice minimises $\sum_{0\ne v\in\Lambda}U(|v|)$ for every
$z_e,z_h,h_1,h_2$.
\end{proposition}
\emph{Proof.} By Theorem~\ref{thm:main} and Tonelli, $\sum_vU(|v|)=\int_0^\infty\mu(\sigma)\,\Theta_\Lambda(\sigma)\,d\sigma$ with
$\Theta_\Lambda(\sigma)=\sum_{v\ne0}e^{-\sigma|v|^2}$, and $\mu\ge0$. For every $\sigma>0$ and every fixed density the triangular lattice minimises
$\Theta_\Lambda(\sigma)$ among lattices (Rankin, Cassels, Ennola, Diananda; see Montgomery, Glasgow Math.\ J.\ \textbf{30} (1988) 75--85;
references from memory). Integrate. \hfill$\square$

This uses no AI-generated input. What remains conditional is the extension from lattices to \emph{all} periodic configurations.

\begin{proposition}[conditional]
\emph{If} the triangular lattice minimises, among periodic configurations of given density in $\R^2$, the energy per particle for every
$f(r^2)$ with $f$ completely monotone (and a convergence condition at $0$ and $\infty$), \emph{then} it does so for the gated kernel
$U$ of Theorem~\ref{thm:main} for every $z_e,z_h,h_1,h_2$, since $U(\sqrt t)$ is completely monotone and decays exponentially.
\end{proposition}

This is a statement about periodic, bulk, classical configurations of the ideal-metal model. It says nothing about open boundaries,
quantum states, or driven-dissipative dynamics, and a hexagonal pattern in an experiment does not single out complete monotonicity
(the triangular lattice is the ground state of the Coulomb, Yukawa, dipolar and Gaussian-mixture potentials of this type, but not of all monotone repulsions, e.g. soft-core ones).

\section{What we ask the reader}

\begin{enumerate}[itemsep=2pt]
\item Is Theorem~\ref{thm:main} known? It looks like the statement that a Dirichlet Green function in a product domain is a
non-negative mixture of 2D heat kernels, applied to a neutral charge combination; a pointer to the literature would settle priority. (The modal expansion of Proposition~\ref{prop:diel} in $K_0$ modes is the classical parallel-plate expansion; what we believe is new is only the neutral-square positivity argument, the lattice proposition and the formalisation.)
\item Is Proposition~\ref{prop:diel} correct as stated (measurability/regularity of $\varepsilon$, endpoints, the weight $\varepsilon\,dz$)?
\item Is the status of the planar universal-optimality statement known to the reader, and is the conditional form above the right way to
state the consequence?
\item Is the Lean statement a faithful encoding of \eqref{eq:images}--\eqref{eq:U}? (The kernel is the combination of six lattice sums
\texttt{Uimg}; the verifier checked it numerically, not by a second formal proof.)
\end{enumerate}

\begin{thebibliography}{9}\small
\bibitem{OpenAI2026} OpenAI, \emph{Universal optimality of the triangular lattice}, \url{https://github.com/openai/math}, family 090 (2026); AI-generated, unrefereed.
\bibitem{CohnKumar2007} H.~Cohn and A.~Kumar, \emph{Universally optimal distribution of points on spheres}, J. Amer. Math. Soc. \textbf{20} (2007), 99--148.
\bibitem{CKMRV2022} H.~Cohn, A.~Kumar, S.~D.~Miller, D.~Radchenko, M.~Viazovska, \emph{Universal optimality of the $E_8$ and Leech lattices and interpolation formulas}, Ann. of Math. \textbf{196} (2022), 983--1082.
\end{thebibliography}
\end{document}
